% ============================================================================
\section{Related Work}
\label{sec:related}

\paragraph{Recurrent-depth and looped transformers.}
Universal Transformers \citep{dehghani2019universal} combine depth-wise weight sharing with adaptive computation.
Reasoning with Latent Thoughts \citep{saunshi2025latentthoughts} and Loop, Think, \& Generalize \citep{kohli2026loopthink} show that looped transformers can trade repeated latent computation for feedforward depth on reasoning-style tasks.
Huginn \citep{geiping2025huginn} scales recurrent-depth language modeling and frames latent reasoning as a third test-time compute axis.
Ouro \citep{zhu2025ouro} pretrains a looped language model with inter-loop normalization and entropy-regularized adaptive exit.
LoopFormer \citep{jeddi2026loopformer} trains budget-conditioned looped models with shortcut modulation.
Elastic Looped Transformer (ELT) \citep{goyal2026elt} studies visual generation with looped transformers and uses intra-loop self-distillation to make intermediate loops usable.
Parcae \citep{prairie2026parcae} studies stable looped language-model scaling through a dynamical-systems view and spectral constraints on looped updates.
Recent mechanistic analyses \citep{blayney2026mechanistic,labovich2026stability} study cyclic fixed points, stages of inference, normalization, and stability in looped models.
Our focus is complementary: we study how the supervised output readout determines whether CE can see recurrent hidden-state scale.

\paragraph{Adaptive computation and early exit.}
Adaptive Computation Time \citep{graves2016adaptive}, PonderNet \citep{banino2021pondernet}, and depth-adaptive Transformers \citep{elbayad2020depth} all motivate spending different amounts of computation on different examples.
Our setting differs because every exit is another application of the same recurrent block, so the hidden state used for an early prediction is also the state that would be fed into later computation.
This makes readout geometry part of the adaptive-compute problem.

\paragraph{Normalization and recurrent scale.}
RMSNorm \citep{zhang2019rmsnorm} explicitly provides re-scaling invariance.
LayerNorm placement \citep{xiong2020layernorm} changes transformer optimization, while NormFormer \citep{shleifer2021normformer} and DeepNet \citep{wang2022deepnet} add normalization or residual scaling to improve fixed-depth pretraining.
TaperNorm \citep{kanavalau2026tapernorm} is the closest output-side comparison: it studies removing normalization while maintaining output scale anchoring and avoiding logit chasing in standard transformers.
Our distinction is the recurrent interface.
Output normalization can stabilize fixed-depth transformers; in a looped model, however, the same hidden state is reused by future loops.
Thus the readout is not only an output calibration choice: it determines whether CE can supervise an active recurrent state variable.
When every CE signal for that recurrent state passes through a scale-invariant interface, scale can become an underconstrained recurrent degree of freedom.
Classical recurrent stability work \citep{pascanu2013difficulty} studies exploding and vanishing dynamics; our results show a compatible lever: make the loss itself sensitive to recurrent scale, or remove scale from the recurrent state.

\paragraph{Concurrent work.}
Several concurrent 2026 works study the stability and supervision of looped models from directions complementary to ours.
On the dynamics side, the Fully Looped Transformer \citep{fu2026fullylooped} attributes training instability at high loop counts to gradient oscillation and residual explosion and responds architecturally, while STARS \citep{yang2026stars} regularizes the Jacobian spectral radius so that latent states approach stable fixed points at large test-time depth.
Both intervene on the recurrent dynamics directly; we instead ask which recurrent variables the supervised readout lets any loss control, and show that scale can drift under dense supervision precisely because every supervised interface hides it.
On the supervision side, RLTT \citep{williams2026rltt} and LoopRPT \citep{tang2026looprpt} find that reinforcement objectives on looped models improve when credit is assigned to intermediate latent states rather than only the final state.
This matches our observation that loss placement determines which parts of the latent trajectory are trained; the visibility question we raise, namely what each per-step signal can see through the readout, applies to those objectives as well.
Finally, concurrent scaling studies \citep{schwethelm2026isodepth,lee2026sparse} quantify when looping is worth its training compute and identify loop boundaries as natural early-exit points, reinforcing the usable-intermediate-exit requirement that motivates our variable-depth evaluation.
